% content/c02-gpt/en.tex — Chapter 2: General probabilistic theories (English, follows de.tex)
%
% Source of the definitions and quoted theorems: Plávala, arXiv:2103.07469v2
%   (source key plavala-2103.07469v2; files basic_v3.tex, what_v4.tex, CT_v3.tex,
%   QT_v3.tex, boxworld_v3.tex, channels_v3.tex, cones_v3.tex).
% Mathematics: evidence/derivations/c02/DERIVATION.md, checked in REVIEW.md beside it.
% Framing: docs/plan/BRIEF.md, section "The physics approach", items 0 to 3
%   (author-stipulated; quoted and arranged only, not extended).
%
% Notation (registry, TR-PLA-02 and TR-PLA-04): states x, y; effects f, g, h.
% DERIVATION.md and IDEAS.md write omega_i, e, pi_i for the same quantities; here they are
% x_i, f, lambda_i.  The coordinates (b, a) of an effect are those of the derivation.

\chapter{General probabilistic theories}\label{c02}

This chapter fixes the framework in which the whole course works. In it a theory consists
of two sets: the states into which a system can be put, and the yes-no questions that can
be asked of it afterwards. To every pair of a state and a question belongs a probability
of the answer yes. We work through three theories completely: the classical bit, the
square bit and the qubit. At the end stands a task that tells the three apart:
distinguishing two preparations with a single measurement.

The mathematics follows the introduction to general probabilistic theories by Plávala
\srcref{plavala-2103.07469v2}{“What are GPTs?” and “State spaces and effect algebras”}.
As there, all state spaces of this chapter are finite-dimensional. That is an assumption
and not a conclusion; it is condition (S5) of \defref{D3}.

\section{What a theory supplies}\label{c02.s1}

The course follows a framing that its author has laid down. This chapter does not justify
it; it uses it. Its first sentences read:
% GAP-TJO: the German version of the three sentences below is a translation of
% docs/plan/BRIEF.md:44-46 and was not written by TJO; the English wording here is the
% brief's (the third sentence with a capital letter and typographic quotation marks).
\begin{equation}\label{c02.eq.rahmen-liefert}\claim{c02.cl.rahmen-liefert}{author-stipulated}
  \begin{aligned}
    &\text{\tpcifenglish{A physical theory supplies data for}{Eine physikalische Theorie liefert Daten für}} \\
    &\text{\tpcifenglish{preparations + dynamics + measurement.}{Präparationen + Dynamik + Messung.}}
  \end{aligned}
\end{equation}
\begin{equation}\label{c02.eq.rahmen-probabilistisch}\claim{c02.cl.rahmen-probabilistisch}{author-stipulated}
  \begin{aligned}
    &\text{\tpcifenglish{All physical theories are probabilistic theories}{Alle physikalischen Theorien sind probabilistische Theorien}} \\
    &\text{\tpcifenglish{(I think of deterministic as a subset of probabilistic).}{(deterministische fasse ich als Teilmenge der probabilistischen auf).}}
  \end{aligned}
\end{equation}
\begin{equation}\label{c02.eq.rahmen-sprache}\claim{c02.cl.rahmen-sprache}{author-stipulated}
  \begin{aligned}
    &\text{\tpcifenglish{We use the language of preparations (= states)}{Wir benutzen die Sprache der Präparationen (= Zustände)}} \\
    &\text{\tpcifenglish{instead of “probability measures”.}{statt der „Wahrscheinlichkeitsmaße“.}}
  \end{aligned}
\end{equation}
Of the three ingredients in \eqref{c02.eq.rahmen-liefert} this chapter treats two: the
preparations and the measurements. Dynamics comes later in the course.

What a preparation is, is fixed by \defref{D1} in Chapter~\ref{c00}: a list of
instructions by which a system is set up for an experiment. After that the system is asked
a yes-no question. The theory assigns to every pair of a preparation and a yes-no question
a number between $0$ and $1$, the probability of the answer yes
\srcref{plavala-2103.07469v2}{“Effect algebra”}.
Figure~\ref{c02.fig.schema} shows one run of such an experiment.

\begin{tpcfigure}{c02.fig.schema}
  \caption{One run of an experiment. The preparation puts the system into a state $x$;
    then a yes-no question is asked. The theory predicts the probability $f(x)$ of the
    answer yes. Throughout the chapter states are drawn solid and as filled dots, effects
    dashed and as squares.}
\end{tpcfigure}

Two preparations that give the same probabilities for all yes-no questions count as the
same state by \defref{D1}. A state is thus what matters about a preparation for the
predictions.

\section{State spaces}\label{c02.s2}

The set of all states of a system is called its state space.

\tpcdef{D3}

In the examples of this chapter $\StSp$ is a subset of $\mathbb{R}^d$ with $d = 1$, $2$ or
$3$. Closed and bounded then means: $\StSp$ contains its boundary and does not extend to
infinity.

Why convex? If $x$ and $y$ are states and $\lambda$ is a number between $0$ and $1$, one
can roll dice: with probability $\lambda$ the state $x$ is prepared, with probability
$1-\lambda$ the state $y$ is prepared. That is a preparation too. Its state is described
by the point $\lambda x + (1-\lambda) y$
\srcref{plavala-2103.07469v2}{“What are GPTs?”}.
Geometrically this point lies on the segment between $x$ and $y$, the closer to $x$ the
larger $\lambda$ is. Convex means: with any two points, $\StSp$ contains the whole segment
joining them.

\tpcdef{D4}

\begin{tpcfigure}{c02.fig.zustandsraeume}
  \demo{c02.demo.explorer}
  \caption{The three state spaces of this chapter: (a) a segment, (b) a square, (c) a
    ball, drawn here with its equator. In each the mixture $z = \tfrac14 x + \tfrac34 y$
    of two states $x$ and $y$ is marked (open circle); it lies on the segment joining
    them, three times as close to $y$ as to $x$. The pure states are the two end points
    $s_0$, $s_1$ in (a), the four corners $s_{00}$, $s_{10}$, $s_{01}$, $s_{11}$ in (b),
    and all points of the surface of the ball in (c). This is verified in
    Section~\ref{c02.s4}.}
\end{tpcfigure}

Pure states are the corners of the state space in the widest sense: points that lie in the
interior of no segment joining two other states. All other states can be mixed from them:

\begin{theorem}[Every state is a mixture of pure states]\label{c02.thm.zerlegung}\claim[proof-omitted]{c02.cl.zerlegung}{source-backed}
  Let $\StSp \subset V$ be a state space in a real vector space $V$ of dimension $d$.
  Then every state $x \in \StSp$ can be written as a convex combination of at most $d+1$
  pure states: there are pure states $y_1, \dots, y_k$ with $k \leq d+1$ and numbers
  $\lambda_1, \dots, \lambda_k \geq 0$ with $\lambda_1 + \dots + \lambda_k = 1$ and
  \begin{equation}\label{c02.eq.zerlegung}
    x = \lambda_1 y_1 + \dots + \lambda_k y_k .
  \end{equation}
\end{theorem}

These are two theorems of convex geometry. Plávala takes them over without proof
\srcref{plavala-2103.07469v2}{“State space”}.
In the square one therefore always manages with at most three corners, in the ball with at
most four points of the surface. In general the decomposition is not unique; the square
will supply the first counterexample in equation~\eqref{c02.eq.quadrat-mitte}.

\section{Effects}\label{c02.s3}

A yes-no question assigns to every state $x$ a number $f(x)$, the probability of the
answer yes. For the theory it is therefore a function $f$ from $\StSp$ to $[0,1]$. Not
every such function qualifies.

To see this we look at the dice preparation from above in two ways. In every run the die
decides what is prepared: with probability $\lambda$ it shows the label “$x$”, and $x$ is
prepared; with probability $1-\lambda$ it shows the label “$y$”, and $y$ is prepared.

First calculation. A run carries the label “$x$” with probability $\lambda$ and then gives
a yes with probability $f(x)$. With probability $1-\lambda$ it carries the label “$y$” and
then gives a yes with probability $f(y)$. Together the probability of yes is
$\lambda f(x) + (1-\lambda) f(y)$.

Second calculation. The state that was prepared is $\lambda x + (1-\lambda) y$. The
probability of yes is $f(\lambda x + (1-\lambda) y)$.

Both calculations describe the same experiment. So for all $x, y \in \StSp$ and all
$\lambda \in [0,1]$ we must have
\begin{equation}\label{c02.eq.affin}\claim{c02.cl.affin}{source-quoted}
  f(\lambda x + (1-\lambda) y) = \lambda f(x) + (1-\lambda) f(y) .
\end{equation}
A function with this property is called \emph{affine}
\srcref{plavala-2103.07469v2}{“Effect algebra”}.

\begin{tpcfigure}{c02.fig.affin}
  \caption{Mixing the states or mixing the probabilities. Horizontally the segment between two states is plotted, vertically the probability of yes. (a) An affine function: its graph is a straight line. In the mixture $z = \tfrac14 x + \tfrac34 y$ it has the value $\tfrac14\cdot\tfrac14 + \tfrac34\cdot\tfrac34 = \tfrac58$, the same mixture of the two end values. (b) The function $r \mapsto r^2$ on $[-1,1]$: the mean of the end values is $1$, the value in the middle is $0$. Such a function cannot describe a yes-no question.}
\end{tpcfigure}

\begin{greybox}[Two assumptions behind equation~\eqref{c02.eq.affin}]
  The derivation uses two assumptions, both of which Plávala states explicitly.
  \begin{enumerate}
    \item Mixtures can be produced by repeating the experiment often and preparing
      sometimes $x$, sometimes $y$.
    \item The result is the same whether one mixes the states or the probabilities. In
      other words: the probability of an answer depends only on the state and not on the
      way in which it was produced.
  \end{enumerate}
  Equation~\eqref{c02.eq.affin} does not follow from the convexity of $\StSp$ alone. On
  the segment $[-1,1]$, for instance, the function $r \mapsto r^2$ has values in $[0,1]$,
  but midway between $-1$ and $1$ it has the value $0$ and not the mean $1$ of the two
  boundary values.
\end{greybox}

\tpcdef{D2}

In coordinates affine functions are easy to recognise\zoomref{c02.v1}.

\begin{lemma}[Affine functions in coordinates]\label{c02.lem.koordinaten}\claim[proof-linked]{c02.cl.koordinaten}{derived}
  Let $\StSp \subset \mathbb{R}^d$ be a state space that is contained in no line ($d = 2$),
  in no plane ($d = 3$) and in general in no affine subspace of smaller dimension; for
  $d = 1$ this means that $\StSp$ has more than one point. A function
  $f\colon \StSp \to \mathbb{R}$ satisfies equation~\eqref{c02.eq.affin} if and only if
  there are a number $b$ and a vector $\vb{a} \in \mathbb{R}^d$ with
  \begin{equation}\label{c02.eq.koordinaten}
    f(\vb{r}) = b + \vb{a}\cdot\vb{r}
    \qquad \text{\tpcifenglish{for all}{für alle}}\ \vb{r} \in \StSp .
  \end{equation}
  The number $b$ and the vector $\vb{a}$ are uniquely determined by $f$.
\end{lemma}

\begin{vertiefung}[id=c02.v1,
    titel={Why affine functions have this form},
    wozu={The core uses Lemma~\ref{c02.lem.koordinaten} without proof. Here it is proved,
      together with the rule for mixtures of more than two states.},
    voraussetzung={Core up to Lemma~\ref{c02.lem.koordinaten}; systems of linear
      equations. The rest of the core does not need this block.},
    zurueck={On a state space in $\mathbb{R}^d$ an effect is fixed by $d+1$ numbers: $b$
      and the components of $\vb{a}$.}]
  Equation~\eqref{c02.eq.affin} speaks of two states. For $k$ states $x_1, \dots, x_k$ and
  weights $\lambda_i \geq 0$ with $\lambda_1 + \dots + \lambda_k = 1$ it implies
  \begin{equation}\label{c02.eq.mehrfach}\claim{c02.cl.mehrfach}{derived}
    f(\lambda_1 x_1 + \dots + \lambda_k x_k) = \lambda_1 f(x_1) + \dots + \lambda_k f(x_k) .
  \end{equation}
  This is shown step by step. For $k = 2$ it is equation~\eqref{c02.eq.affin}. If
  $\lambda_k < 1$, then
  $\lambda_1 x_1 + \dots + \lambda_k x_k = (1-\lambda_k)\, y + \lambda_k x_k$ with the state
  $y = \bigl(\lambda_1 x_1 + \dots + \lambda_{k-1} x_{k-1}\bigr)/(1-\lambda_k)$, which is a
  convex combination of $k-1$ states. Equation~\eqref{c02.eq.affin} reduces the claim for
  $k$ to the claim for $k-1$. For $\lambda_k = 1$ there is nothing to show.

  \begin{proof}[Proof of Lemma~\ref{c02.lem.koordinaten}]
    Every function of the form \eqref{c02.eq.koordinaten} is affine, because
    \begin{equation*}
      b + \vb{a}\cdot(\lambda\vb{r} + (1-\lambda)\vb{s})
      = \lambda(b + \vb{a}\cdot\vb{r}) + (1-\lambda)(b + \vb{a}\cdot\vb{s}) .
    \end{equation*}

    Conversely let $f$ be affine. By assumption $\StSp$ contains points
    $\vb{v}_0, \dots, \vb{v}_d$ that lie in no affine subspace of smaller dimension. Then
    the system of linear equations $b + \vb{a}\cdot\vb{v}_i = f(\vb{v}_i)$,
    $i = 0, \dots, d$, for the $d+1$ unknowns $b$ and $\vb{a}$ has exactly one solution. We
    put $g(\vb{r}) = b + \vb{a}\cdot\vb{r}$ and show $f = g$ on all of $\StSp$.

    Every point $\vb{r} \in \mathbb{R}^d$ can be written in exactly one way as
    $\vb{r} = \alpha_0 \vb{v}_0 + \dots + \alpha_d \vb{v}_d$ with
    $\alpha_0 + \dots + \alpha_d = 1$; the $\alpha_i$ may be negative. We move the terms
    with negative $\alpha_i$ to the left-hand side. Let $N$ be the set of these indices and
    $S = 1 + \sum_{i \in N} \abs{\alpha_i}$; the remaining $\alpha_i$ then also add up to
    $S$. Division by $S$ gives
    \begin{equation*}
      \frac{1}{S}\,\vb{r} + \sum_{i \in N} \frac{\abs{\alpha_i}}{S}\,\vb{v}_i
      = \sum_{i \notin N} \frac{\alpha_i}{S}\,\vb{v}_i .
    \end{equation*}
    For $\vb{r} \in \StSp$ both sides are convex combinations of states. Both $f$ and $g$
    satisfy equation~\eqref{c02.eq.mehrfach}, and they agree in all $\vb{v}_i$. Applying
    both to this equation and subtracting the results leaves
    $\bigl(f(\vb{r}) - g(\vb{r})\bigr)/S = 0$.
  \end{proof}
\end{vertiefung}

By equation~\eqref{c02.eq.koordinaten}, $f(\vb{r}) = b + \vb{a}\cdot\vb{r}$. If
$\vb{a} \neq 0$, that is, if $f$ is not constant, the points of $\mathbb{R}^d$ in which
$b + \vb{a}\cdot\vb{r}$ has the same value form parallel lines ($d = 2$) or parallel
planes ($d = 3$), perpendicular to $\vb{a}$; the level sets of $f$ are their intersections
with $\StSp$. If $\vb{a} = 0$, then $f$ is constant, and all of $\StSp$ is a single level
set. An effect is an affine function whose values stay between $0$ and $1$ on all of $\StSp$
(Figure~\ref{c02.fig.niveaulinien}). Affine is not the same as linear: the constant part
$b$ belongs to it.

\begin{tpcfigure}{c02.fig.niveaulinien}
  \caption{An effect is an affine function with values in $[0,1]$. (a) On the segment its
    graph is a straight line; it is fixed by the values $q_0$ and $q_1$ in the end points,
    here $\tfrac34$ and $\tfrac14$. In the state $x$, a quarter of the way from $s_0$ to
    $s_1$, one reads off $f(x) = \tfrac58$. (b) On the square the level lines of the effect
    $f(\vb{r}) = \tfrac12 + \tfrac14 r_x + \tfrac14 r_y$ are parallel lines perpendicular
    to $\vb{a} = (\tfrac14, \tfrac14)$. The value $0$ is taken only in the lower left
    corner, the value $1$ only in the upper right corner; in the state
    $x = (\tfrac12, 0)$ we have $f(x) = \tfrac58$.}
\end{tpcfigure}

\tpcdef{D58}

For every effect $f$ the function $\UnitEff{\StSp} - f$ is an effect as well, because it is
affine and has values in $[0,1]$:
\begin{equation}\label{c02.eq.komplement}\claim{c02.cl.komplement}{source-backed}
  \UnitEff{\StSp} - f \in \EffAlg{\StSp},
  \qquad f + (\UnitEff{\StSp} - f) = \UnitEff{\StSp} .
\end{equation}
The effect $\UnitEff{\StSp} - f$ belongs to the answer no: if $f(x)$ is the probability of
yes, then $1 - f(x)$ is that of no. A yes-no question is thus the same as a measurement
with two outcomes, described by the pair $f$, $\UnitEff{\StSp} - f$. Two effects exist in
every theory: $\UnitEff{\StSp}$, the question with the certain answer yes, and the zero
function $0$, the question with the certain answer no.

Effects can be mixed too. One rolls dice to decide which of two questions is asked: with
probability $\lambda$ the question $f$, otherwise the question $g$. The probability of yes
is then $\lambda f(x) + (1-\lambda) g(x)$, and
\begin{equation}\label{c02.eq.effekt-mischung}\claim{c02.cl.effekt-mischung}{derived}
  \lambda f + (1-\lambda) g \in \EffAlg{\StSp}
  \qquad \text{\tpcifenglish{for all}{für alle}}\ f, g \in \EffAlg{\StSp},\ \lambda \in [0,1] ,
\end{equation}
because this function is again affine and has values in $[0,1]$. The set
$\EffAlg{\StSp}$ is therefore itself convex, and it has corners:

\tpcdef{D74}

\section{Three theories}\label{c02.s4}

\subsection{The classical bit}\label{c02.s4.1}

\tpcdef{D11}

The classical bit is the classical theory with $n = 2$. A state is a probability
distribution on two possibilities, as in Chapter~\ref{c01}:
\begin{equation}\label{c02.eq.bit}\claim{c02.cl.bit-zustaende}{source-backed}
  \begin{aligned}
    \Simplex{2} &= \bigl\{ (\pw_0, \pw_1) \in \mathbb{R}^2 : \pw_0 \geq 0,\ \pw_1 \geq 0,\ \pw_0 + \pw_1 = 1 \bigr\}, \\
    s_0 &= (1,0), \qquad s_1 = (0,1) .
  \end{aligned}
\end{equation}
This is the segment between $s_0$ and $s_1$ in Figure~\ref{c02.fig.zustandsraeume}(a);
every state is $x = \pw_0 s_0 + \pw_1 s_1$, in exactly one way. The pure states are the two
end points $s_0$ and $s_1$
\srcref{plavala-2103.07469v2}{“Example: classical theory”}.

\begin{proposition}[Effects of the classical bit]\label{c02.prop.bit}\claim[proof-included]{c02.cl.bit-effekte}{derived}
  Every effect $f$ on $\Simplex{2}$ is determined by its values $q_0 = f(s_0)$ and
  $q_1 = f(s_1)$ in the pure states:
  \begin{equation}\label{c02.eq.bit-effekt}
    f(\pw_0, \pw_1) = q_0\, \pw_0 + q_1\, \pw_1 .
  \end{equation}
  Conversely this is an effect for every pair $(q_0, q_1)$ with $0 \leq q_0 \leq 1$ and
  $0 \leq q_1 \leq 1$. The set of effects is thus the unit square
  $\EffAlg{\Simplex{2}} = [0,1]^2$. Its extreme effects are the four corners
  $(0,0)$, $(1,0)$, $(0,1)$ and $(1,1)$.
\end{proposition}

\begin{tpcfigure}{c02.fig.effektmengen}
  \caption{The effect sets of two theories. (a) Classical bit: the unit square in the
    coordinates $q_0 = f(s_0)$, $q_1 = f(s_1)$. The point $f$ is the effect of
    Figure~\ref{c02.fig.niveaulinien}(a). (b) Square bit, anticipating
    Proposition~\ref{c02.prop.quadrat}: a double pyramid; $b$ is plotted upwards, $a_x$
    and $a_y$ horizontally. The extreme effects are marked as squares, hidden edges are
    dotted. The bit has four extreme effects for two pure states, the square bit six for
    four pure states.}
\end{tpcfigure}

\begin{proof}
  Equation~\eqref{c02.eq.bit-effekt} is equation~\eqref{c02.eq.affin} for
  $x = \pw_0 s_0 + \pw_1 s_1$. Conversely the right-hand side of
  \eqref{c02.eq.bit-effekt} is affine, and its value is a weighted mean of $q_0$ and
  $q_1$; it lies in $[0,1]$ for all states if and only if $q_0$ and $q_1$ lie there.

  On the corners: if one of the two numbers, say $q_0$, lies strictly between $0$ and $1$,
  it can be increased and decreased by a small amount $\varepsilon$ without leaving the
  square; $(q_0, q_1)$ is then midway between $(q_0 + \varepsilon, q_1)$ and
  $(q_0 - \varepsilon, q_1)$ and not an extreme effect. If on the other hand $q_0$ and
  $q_1$ are both $0$ or $1$ and $f = \lambda g + (1-\lambda) h$ with effects $g$, $h$ and
  $0 < \lambda < 1$, then $g$ and $h$ must have the same value as $f$ in each of the two
  coordinates: a mean of two numbers from $[0,1]$ is $0$ only if both are $0$, and $1$
  only if both are $1$.
\end{proof}

The four extreme effects have a simple meaning. $(1,0)$ asks: “Is it $s_0$?” $(0,1)$
asks: “Is it $s_1$?” $(1,1)$ is the unit effect and $(0,0)$ the zero function. In the pure
states their answers are certain: the probabilities are $0$ or $1$. The other effects are
mixtures of these four, that is, questions in which dice are rolled in addition.
Figure~\ref{c02.fig.effektmengen}(a) shows the set of effects.

For comparison with the other two theories we put the origin at the middle of the segment.
With the coordinate $r = \pw_0 - \pw_1 \in [-1,1]$ we have $\pw_0 = \tfrac12(1+r)$ and
$\pw_1 = \tfrac12(1-r)$, and equation~\eqref{c02.eq.bit-effekt} becomes
\begin{equation}\label{c02.eq.bit-zentriert}\claim{c02.cl.bit-zentriert}{derived}
  \begin{aligned}
    f(r) &= b + a\, r, \qquad b = \tfrac12(q_0 + q_1), \quad a = \tfrac12(q_0 - q_1), \\
    \abs{a} &\leq \min(b,\, 1-b) .
  \end{aligned}
\end{equation}
The inequality says the same as $0 \leq q_0 \leq 1$ and $0 \leq q_1 \leq 1$, because
$q_0 = b + a$ and $q_1 = b - a$ are the largest and the smallest value of $f$ if
$a \geq 0$, and the other way round if $a < 0$.

\subsection{The square bit}\label{c02.s4.2}

The second theory has a square as its state space. Plávala introduces it under the name
“boxworld”: there a state describes a box with one input bit and one output bit, and the
four corners $s_{ij}$ are the boxes that give the output $i$ on input $0$ and the output
$j$ on input $1$
\srcref{plavala-2103.07469v2}{“Example: boxworld theory”}. In coordinates whose origin is
at the centre,
\begin{equation}\label{c02.eq.quadrat}\claim{c02.cl.quadrat-zustaende}{source-backed}
  \begin{aligned}
    \StSp_{\mathrm{Q}} &= [-1,1]^2 = \bigl\{ (r_x, r_y) : \abs{r_x} \leq 1,\ \abs{r_y} \leq 1 \bigr\}, \\
    s_{ij} &= (2i-1,\ 2j-1), \qquad i, j \in \{0,1\} .
  \end{aligned}
\end{equation}
The pure states are the four corners (Figure~\ref{c02.fig.zustandsraeume}(b)). Unlike for
the bit, the decomposition into pure states is not unique. The centre of the square is a
mixture of two corners in two ways:
\begin{equation}\label{c02.eq.quadrat-mitte}\claim{c02.cl.quadrat-mitte}{source-backed}
  \tfrac12\bigl(s_{00} + s_{11}\bigr) = \tfrac12\bigl(s_{10} + s_{01}\bigr) = (0,0) .
\end{equation}
In a simplex this would not be possible: there the weights with which a state is mixed
from the pure states are uniquely determined. By \defref{D11} the square bit is therefore
not a classical theory.

\begin{tpcfigure}{c02.fig.quadrat}
  \caption{The square bit. (a) The centre of the square lies on both diagonals; it is
    midway between $s_{00}$ and $s_{11}$ and also midway between $s_{10}$ and $s_{01}$, as
    in equation~\eqref{c02.eq.quadrat-mitte}. (b) The level lines of the extreme effect
    $f_x^{+}$ of Proposition~\ref{c02.prop.quadrat}: it has the value $1$ on the right
    edge and the value $0$ on the left edge.}
\end{tpcfigure}

\begin{proposition}[Effects of the square bit]\label{c02.prop.quadrat}\claim[proof-included]{c02.cl.quadrat-effekte}{derived}
  The function $f(\vb{r}) = b + a_x r_x + a_y r_y$ is an effect on $\StSp_{\mathrm{Q}}$ if
  and only if
  \begin{equation}\label{c02.eq.quadrat-effekt}
    \abs{a_x} + \abs{a_y} \leq \min(b,\, 1-b) .
  \end{equation}
  If an effect is written as a point $(b, a_x, a_y)$, then
  $\EffAlg{\StSp_{\mathrm{Q}}}$ is a double pyramid with six corners. These six corners are
  the extreme effects:
  \begin{equation}\label{c02.eq.quadrat-extremal}
    \begin{aligned}
      0 &= (0,0,0), & f_x^{\pm} &= \bigl(\tfrac12, \pm\tfrac12, 0\bigr), \\
      \UnitEff{\StSp} &= (1,0,0), & f_y^{\pm} &= \bigl(\tfrac12, 0, \pm\tfrac12\bigr) .
    \end{aligned}
  \end{equation}
\end{proposition}
\begin{proof}
  The function $f$ takes its smallest value in the corner in which $r_x$ has the sign
  opposite to $a_x$ and $r_y$ the sign opposite to $a_y$; it is
  $b - \abs{a_x} - \abs{a_y}$. Correspondingly the largest value is
  $b + \abs{a_x} + \abs{a_y}$. Both lie in $[0,1]$ if and only if
  \eqref{c02.eq.quadrat-effekt} holds.

  Every effect is a mixture of the six points in \eqref{c02.eq.quadrat-extremal}. Indeed,
  with $q = \abs{a_x} + \abs{a_y}$,
  \begin{equation*}
    (b, a_x, a_y) = (b - q)\,\UnitEff{\StSp} + 2\abs{a_x}\, f_x^{\pm} + 2\abs{a_y}\, f_y^{\pm} + (1 - b - q)\cdot 0 ,
  \end{equation*}
  where for $f_x^{\pm}$ the sign of $a_x$ and for $f_y^{\pm}$ that of $a_y$ is to be taken;
  if $a_x = 0$ or $a_y = 0$, the term concerned has weight $0$ and is omitted.
  By \eqref{c02.eq.quadrat-effekt} the four weights are non-negative, and their sum is
  $1$. An extreme effect must therefore be one of the six points.

  That all six are extreme is seen at the places where they take the values $0$ and $1$.
  Let $f_x^{+} = \lambda g + (1-\lambda) h$ with effects $g$, $h$ and $0 < \lambda < 1$.
  The effect $f_x^{+}(\vb{r}) = \tfrac12(1 + r_x)$ has the value $1$ in the corners $(1,1)$
  and $(1,-1)$ and the value $0$ in $(-1,1)$. As in the proof of
  Proposition~\ref{c02.prop.bit}, $g$ and $h$ must have the same values there. The three
  equations $b + a_x + a_y = 1$, $b + a_x - a_y = 1$, $b - a_x + a_y = 0$ have only the
  solution $(b, a_x, a_y) = (\tfrac12, \tfrac12, 0)$; hence $g = h = f_x^{+}$. For
  $f_x^{-}$, $f_y^{+}$, $f_y^{-}$ the argument is the same. The zero function is extreme
  because a mean of two functions with values in $[0,1]$ is $0$ everywhere only if both
  are $0$ everywhere; for $\UnitEff{\StSp}$ the same holds with the value $1$.
\end{proof}

Here too the extreme effects have a simple meaning. The effect
$f_x^{+} = \tfrac12(1 + r_x)$ has the value $i$ in the corner $s_{ij}$: it asks whether the
box gives the output $1$ on input $0$. Correspondingly $f_y^{+} = \tfrac12(1 + r_y)$ has
the value $j$ in $s_{ij}$ and asks for the output on input $1$. $f_x^{-}$ and $f_y^{-}$ are
the complements, $f_x^{-} = \UnitEff{\StSp} - f_x^{+}$ and
$f_y^{-} = \UnitEff{\StSp} - f_y^{+}$. Plávala gives the same set as the convex hull of
these six effects; it is drawn in Figure~\ref{c02.fig.effektmengen}(b).

\subsection{The qubit}\label{c02.s4.3}

The third theory is the smallest nontrivial system of quantum theory. Its states are matrices:
complex $2\times 2$ matrices $\rho$ with three properties,
\begin{equation}\label{c02.eq.qubit}\claim{c02.cl.qubit-zustaende}{source-backed}
  \rho\ \text{\tpcifenglish{Hermitian}{hermitesch}}, \qquad \rho \geq 0, \qquad \tr(\rho) = 1 .
\end{equation}
Hermitian means: the diagonal entries are real, and the other two entries are complex
conjugates of each other. For a Hermitian $2\times 2$ matrix the condition $\rho \geq 0$
means: its two eigenvalues are non-negative. Such matrices are called density matrices
\srcref{plavala-2103.07469v2}{“Example: quantum theory”}.
To survey them we use the unit matrix and the three Pauli matrices
\srcref{tpc-script-v0.2}{Lecture 17}:
\begin{equation}\label{c02.eq.pauli}\claim{c02.cl.pauli}{source-backed}
  \Id = \begin{pmatrix} 1 & 0 \\ 0 & 1 \end{pmatrix}, \quad
  \pauli{x} = \begin{pmatrix} 0 & 1 \\ 1 & 0 \end{pmatrix}, \quad
  \pauli{y} = \begin{pmatrix} 0 & -i \\ i & 0 \end{pmatrix}, \quad
  \pauli{z} = \begin{pmatrix} 1 & 0 \\ 0 & -1 \end{pmatrix} .
\end{equation}

The next proposition says that these matrices form a ball, as shown in Figure~\ref{c02.fig.blochzustaende}.

\begin{tpcfigure}{c02.fig.blochzustaende}
  \caption{The states of the qubit as points of the unit ball. The points $\vb{n}$, $-\vb{n}$, $\vb{m}$, $-\vb{m}$ of the surface are pure states. The state $\vb{r} = \tfrac12\vb{n}$ in the interior is the mixture $\tfrac34\vb{n} + \tfrac14(-\vb{n})$. The centre is midway between $\vb{n}$ and $-\vb{n}$ and also midway between $\vb{m}$ and $-\vb{m}$.}
\end{tpcfigure}

\begin{proposition}[The qubit state ball]\label{c02.prop.bloch}\claim[proof-included]{c02.cl.bloch}{derived}
  \begin{enumerate}
    \item Every Hermitian $2\times 2$ matrix with trace $1$ can be written with exactly
      one vector $\vb{r} = (r_x, r_y, r_z) \in \mathbb{R}^3$ as
      \begin{equation}\label{c02.eq.bloch}
        \rho(\vb{r}) = \tfrac12\bigl(\Id + r_x \pauli{x} + r_y \pauli{y} + r_z \pauli{z}\bigr)
        = \frac12 \begin{pmatrix} 1 + r_z & r_x - i r_y \\ r_x + i r_y & 1 - r_z \end{pmatrix} .
      \end{equation}
    \item The eigenvalues of $\rho(\vb{r})$ are $\tfrac12(1 + \norm{\vb{r}})$ and
      $\tfrac12(1 - \norm{\vb{r}})$. Hence $\rho(\vb{r})$ is a state if and only if
      $\norm{\vb{r}} \leq 1$.
    \item Mixtures stay mixtures:
      $\rho(\lambda \vb{r} + (1-\lambda)\vb{s}) = \lambda \rho(\vb{r}) + (1-\lambda)\rho(\vb{s})$.
    \item The state $\rho(\vb{r})$ is pure if and only if $\norm{\vb{r}} = 1$.
  \end{enumerate}
\end{proposition}
\begin{proof}
  (1) A Hermitian matrix has real diagonal entries; if its trace is $1$, they can be
  written as $\tfrac12(1 + r_z)$ and $\tfrac12(1 - r_z)$, where $r_z$ is their difference.
  The lower left entry is a complex number, which we call $\tfrac12(r_x + i r_y)$; the
  upper right entry is then the conjugate number $\tfrac12(r_x - i r_y)$.

  (2) The eigenvalues $\kappa$ are the zeros of
  \begin{equation*}
    \det\bigl(\kappa\, \Id - \rho(\vb{r})\bigr)
    = \bigl(\kappa - \tfrac12(1 + r_z)\bigr)\bigl(\kappa - \tfrac12(1 - r_z)\bigr) - \tfrac14\bigl(r_x^2 + r_y^2\bigr)
    = \bigl(\kappa - \tfrac12\bigr)^2 - \tfrac14 \norm{\vb{r}}^2 .
  \end{equation*}
  Both are non-negative if and only if $\norm{\vb{r}} \leq 1$.

  (3) The right-hand side of \eqref{c02.eq.bloch} is a constant plus an expression that is
  linear in $\vb{r}$.

  (4) Let $\norm{\vb{r}} = 1$ and $\vb{r} = \lambda \vb{s} + (1-\lambda)\vb{t}$ with
  $\norm{\vb{s}} \leq 1$, $\norm{\vb{t}} \leq 1$ and $0 < \lambda < 1$. Then
  $1 = \vb{r}\cdot\vb{r} = \lambda\, \vb{r}\cdot\vb{s} + (1-\lambda)\, \vb{r}\cdot\vb{t}$.
  A scalar product is at most the product of the lengths, here at most $1$, and it equals
  $1$ only if the two vectors are equal. Hence $\vb{s} = \vb{t} = \vb{r}$, and
  $\rho(\vb{r})$ is pure. If on the other hand $q = \norm{\vb{r}} < 1$, we choose a unit
  vector $\vb{n}$ with $\vb{r} = q\,\vb{n}$ and write
  $\vb{r} = \tfrac12(1+q)\,\vb{n} + \tfrac12(1-q)\,(-\vb{n})$. This is a mixture of two
  different states with weights strictly between $0$ and $1$.
\end{proof}

Because of (3) we may identify the state $\rho(\vb{r})$ with the point $\vb{r}$. The state
space of the qubit is then the unit ball
$\StSp_{\mathrm{B}} = \{ \vb{r} \in \mathbb{R}^3 : \norm{\vb{r}} \leq 1 \}$. We call it
the qubit state ball; its pure states are the points of the surface
(Figure~\ref{c02.fig.zustandsraeume}(c)). Here too the decomposition into pure states is
not unique: the centre is midway between $\vb{n}$ and $-\vb{n}$ for every unit vector
$\vb{n}$.

The effects are found as for the square, only with a different length for $\vb{a}$. They
too can be written as matrices. Figure~\ref{c02.fig.blochkugel} shows in advance what an
extreme effect of the ball looks like.

\begin{tpcfigure}{c02.fig.blochkugel}
  \caption{The qubit state ball with the extreme effect
    $f_{\vb{n}}(\vb{r}) = \tfrac12(1 + \vb{n}\cdot\vb{r})$ of
    Proposition~\ref{c02.prop.qubit-effekte} for $\vb{n} = (0,0,1)$, that is,
    $f_{\vb{n}}(\vb{r}) = \tfrac12(1 + r_z)$. The points of equal probability form
    horizontal planes. The plane $f = 1$ touches the ball in the pure state $\vb{n}$
    (top), the plane $f = 0$ in the pure state $-\vb{n}$ (bottom). The state $x$ lies at
    height $r_z = \tfrac12$; there $f = \tfrac34$. For an effect that is not
    constant and not extreme, at most one of the two planes $f = 0$ and $f = 1$ touches the
    ball; one of them may still touch it. A constant effect has the same value everywhere;
    its level sets are not planes.}
\end{tpcfigure}

\begin{proposition}[Effects of the qubit]\label{c02.prop.qubit-effekte}\claim[proof-included]{c02.cl.qubit-effekte}{derived}
  \begin{enumerate}
    \item The function $f(\vb{r}) = b + \vb{a}\cdot\vb{r}$ is an effect on
      $\StSp_{\mathrm{B}}$ if and only if
      \begin{equation}\label{c02.eq.qubit-effekt}
        \norm{\vb{a}} \leq \min(b,\, 1-b) .
      \end{equation}
    \item With the Hermitian matrix
      $\op{\EffOp} = b\, \Id + a_x \pauli{x} + a_y \pauli{y} + a_z \pauli{z}$,
      \begin{equation}\label{c02.eq.spurformel}
        f(\vb{r}) = b + \vb{a}\cdot\vb{r} = \tr\bigl(\rho(\vb{r})\, \op{\EffOp}\bigr) .
      \end{equation}
      The eigenvalues of $\op{\EffOp}$ are $b + \norm{\vb{a}}$ and $b - \norm{\vb{a}}$;
      condition~\eqref{c02.eq.qubit-effekt} says that both lie in $[0,1]$.
    \item The extreme effects are $0$, $\UnitEff{\StSp}$ and, for every unit vector
      $\vb{n}$, the effect
      \begin{equation}\label{c02.eq.qubit-extremal}
        f_{\vb{n}}(\vb{r}) = \tfrac12\bigl(1 + \vb{n}\cdot\vb{r}\bigr) .
      \end{equation}
  \end{enumerate}
\end{proposition}
\begin{proof}
  (1) On the ball $\vb{a}\cdot\vb{r}$ is at most $\norm{\vb{a}}$ and at least
  $-\norm{\vb{a}}$; both values are taken, for $\vb{a} \neq 0$ in the points
  $\pm\vb{a}/\norm{\vb{a}}$. The smallest value of $f$ is therefore $b - \norm{\vb{a}}$,
  the largest $b + \norm{\vb{a}}$.

  (2) Multiplying out the matrices in \eqref{c02.eq.pauli} one finds: every Pauli matrix
  has trace $0$, its square is $\Id$, and the product of two different Pauli matrices has
  trace $0$. So $\tr(\Id) = 2$, the trace of the square of a Pauli matrix is $2$, and in
  the product $\rho(\vb{r})\,\op{\EffOp}$ only the terms with $\Id\cdot\Id$ and with two equal
  Pauli matrices contribute to the trace:
  $\tr(\rho(\vb{r})\,\op{\EffOp}) = \tfrac12\bigl(2b + 2 a_x r_x + 2 a_y r_y + 2 a_z r_z\bigr)$.
  The eigenvalues of $\op{\EffOp}$ are obtained as in Proposition~\ref{c02.prop.bloch}(2).

  (3) Let $q = \norm{\vb{a}}$. For $q = 0$,
  $(b, 0) = b\, \UnitEff{\StSp} + (1 - b)\cdot 0$, a mixture of $\UnitEff{\StSp}$ and the
  zero function. For $q > 0$ we put $\vb{n} = \vb{a}/q$; then
  \begin{equation*}
    (b, \vb{a}) = (b - q)\, \UnitEff{\StSp} + 2q\, \bigl(\tfrac12, \tfrac12\vb{n}\bigr) + (1 - b - q)\cdot 0 ,
  \end{equation*}
  a mixture with non-negative weights of sum $1$. Extreme effects can therefore exist only
  among those named. The effect $f_{\vb{n}}$ has the value $1$ in the point $\vb{n}$ and
  the value $0$ in the point $-\vb{n}$. If $f_{\vb{n}} = \lambda g + (1-\lambda) h$ with
  effects $g$, $h$ and $0 < \lambda < 1$, then $g$ and $h$ have the same values there. For
  $g = (b, \vb{a})$ this means $b + \vb{a}\cdot\vb{n} = 1$ and $b - \vb{a}\cdot\vb{n} = 0$,
  so $b = \tfrac12$ and $\vb{a}\cdot\vb{n} = \tfrac12$. By \eqref{c02.eq.qubit-effekt} we
  have $\norm{\vb{a}} \leq \tfrac12$, and that leaves only $\vb{a} = \tfrac12\vb{n}$. Hence
  $g = h = f_{\vb{n}}$. For $0$ and $\UnitEff{\StSp}$ the argument of
  Proposition~\ref{c02.prop.quadrat} applies.
\end{proof}

The extreme effect $f_{\vb{n}}$ asks: “Is it the pure state $\vb{n}$?” In the state
$\vb{n}$ the answer is certainly yes, in the opposite state $-\vb{n}$ certainly no, and in
every other state it is random. Its matrix is
$\tfrac12(\Id + n_x\pauli{x} + n_y\pauli{y} + n_z\pauli{z}) = \rho(\vb{n})$, the density
matrix of this pure state. In Figure~\ref{c02.fig.blochkugel} it is drawn as a family of
parallel planes.

\subsection{What distinguishes the three}\label{c02.s4.4}

In coordinates with the origin at the centre the three theories look almost alike. Each
time an effect is a function $b + \vb{a}\cdot\vb{r}$ with the condition that a length of
$\vb{a}$ is at most $\min(b, 1-b)$. Only the length is measured differently: for the bit as
the modulus $\abs{a}$ (equation~\eqref{c02.eq.bit-zentriert}), for the square bit as the
sum $\abs{a_x} + \abs{a_y}$ (equation~\eqref{c02.eq.quadrat-effekt}), for the qubit as the
ordinary length $\norm{\vb{a}}$ (equation~\eqref{c02.eq.qubit-effekt}).
Table~\ref{c02.tab.vergleich} collects what follows from this.

\begin{tpctable}{c02.tab.vergleich}
  \caption{The three theories compared. The entries are the results of
    Propositions~\ref{c02.prop.bit}, \ref{c02.prop.quadrat}, \ref{c02.prop.bloch} and
    \ref{c02.prop.qubit-effekte} and of equations~\eqref{c02.eq.bit} and
    \eqref{c02.eq.quadrat-mitte}.}
  \begin{tabular}{llll}
    \toprule
      & bit & square bit & qubit \\
    \midrule
    state space & segment & square & ball \\
    pure states & two & four & surface \\
    extreme effects & four & six & $0$, $\UnitEff{\StSp}$, all $f_{\vb{n}}$ \\
    decomposition unique & yes & no & no \\
    classical (\defref{D11}) & yes & no & no \\
    \bottomrule
  \end{tabular}
\end{tpctable}

That the qubit is not a classical theory follows, as for the square, from the ambiguous
decomposition. How the square bit and the qubit differ from each other is shown by the
task in Section~\ref{c02.s6}.

\section{Measurements and the no-restriction hypothesis}\label{c02.s5}

A yes-no question has two possible outcomes. For measurements with more outcomes one needs
a further notion.

\tpcdef{D8}

\tpcdef{D9}

In general, channels describe how states of one system become states of another. This
chapter needs them only for \defref{D9}, and of that only the second half: a measurement
with $n$ outcomes is a list of effects $f_1, \dots, f_n$ whose sum is the unit effect. In
the state $x$ the $i$-th outcome occurs with probability $f_i(x)$, and these numbers form a
probability distribution:
\begin{equation}\label{c02.eq.messung}\claim{c02.cl.messung}{derived}
  f_i(x) \geq 0 \quad (i = 1, \dots, n),
  \qquad f_1(x) + \dots + f_n(x) = \UnitEff{\StSp}(x) = 1 .
\end{equation}
For $n = 2$ this is the yes-no question with $f_1 = f$ and $f_2 = \UnitEff{\StSp} - f$ of
equation~\eqref{c02.eq.komplement}
\srcref{plavala-2103.07469v2}{“Measurements”}.

\begin{tpcfigure}{c02.fig.messung}
  \caption{A measurement with three outcomes on the classical bit. The effects are $f_1 = (\tfrac12, 0)$, $f_2 = (\tfrac12, \tfrac12)$ and $f_3 = (0, \tfrac12)$ in the coordinates $(q_0, q_1)$; their sum is $(1,1) = \UnitEff{\StSp}$. Above every state the three probabilities are stacked; together they fill the height $1$. In the state $x = (\tfrac34, \tfrac14)$ they are $\tfrac38$, $\tfrac12$ and $\tfrac18$.}
\end{tpcfigure}

So far an effect was a function with certain properties. Whether an experiment can be
given for every such function is a question of its own.

\tpcdef{D12}

% GAP-TJO: whether the no-restriction hypothesis is part of postulate P3 in the course or
% a separate assumption is not said in docs/plan/BRIEF.md (BRIEF.md:53).  The text below
% names it only as a hypothesis of this chapter.
This chapter assumes the hypothesis everywhere: in Section~\ref{c02.s6} we optimise over
all effects in $\EffAlg{\StSp}$. If a theory admits only part of the effects, the optimum
computed there is an upper bound and need not be attainable. Moreover the hypothesis
speaks only of yes-no questions. That every list of effects with sum $\UnitEff{\StSp}$ can
also be carried out as a measurement would be a further assumption; the discrimination
task of Section~\ref{c02.s6} uses only measurements with two outcomes and does not need
it\zoomref{c02.v2}. The measurement with three outcomes in
Figure~\ref{c02.fig.messung} is an example for \defref{D9} and does not occur in the
task.

\begin{vertiefung}[id=c02.v2,
    titel={Two cones and the plane at height one},
    wozu={States and effects live in two different vector spaces that are dual to each
      other. Here both spaces and their cones are introduced and computed for the three
      theories.},
    voraussetzung={Core up to Section~\ref{c02.s5}; vector spaces, linear maps. The core
      does not need this block; the blocks~\ref{c02.v3} and \ref{c02.v5} build on it.},
    zurueck={A state is a point of a cone at height one; an effect is a point of another
      cone between $0$ and the unit effect.}]
  \tpcdef{D57}

  \tpcdef{D59}

  The cone $\AffPos{\StSp}$ consists of functions on states; the effects lie in it. The
  cone $\AffDualPos{\StSp}$ consists of functionals on these functions; the states lie in
  it. The two must not be confused. Where the states lie in their cone is said by a theorem
  of Plávala
  \srcref{plavala-2103.07469v2}{“Duality between the state space and the effect algebra”}:

  \begin{theorem}[States at height one]\label{c02.thm.hoehe}\claim[proof-omitted]{c02.cl.hoehe}{source-quoted}
    For every state space $\StSp$ we have, up to an isomorphism,
    \begin{equation}\label{c02.eq.hoehe}
      \StSp = \bigl\{ \varphi \in \AffDualPos{\StSp} : \pair{\varphi}{\UnitEff{\StSp}} = 1 \bigr\} .
    \end{equation}
  \end{theorem}

  The isomorphism is the evaluation of \defref{D59}: the state $x$ becomes the functional
  $f \mapsto f(x)$. The condition $\pair{\varphi}{\UnitEff{\StSp}} = 1$ is the
  normalisation; it cuts a plane out of the cone, the “height one”.

  \begin{lemma}[The cone of unnormalised states]\label{c02.lem.basis}\claim[proof-included]{c02.cl.basis}{derived}
    Every $\psi \in \AffDualPos{\StSp}$ with $\psi \neq 0$ can be written in exactly one
    way as $\psi = t\, x$ with a number $t > 0$ and a state $x \in \StSp$. Here
    $t = \pair{\psi}{\UnitEff{\StSp}}$.
  \end{lemma}
  \begin{proof}
    For every $f \in \Aff{\StSp}$ there is a number $M \geq 0$ with
    $-M \leq f(x) \leq M$ for all $x \in \StSp$, because by
    Lemma~\ref{c02.lem.koordinaten} $f$ is continuous and $\StSp$ is closed and bounded.
    The functions $M\,\UnitEff{\StSp} - f$ and $M\,\UnitEff{\StSp} + f$ then lie in
    $\AffPos{\StSp}$, and for $\psi \in \AffDualPos{\StSp}$ it follows that
    $-M\, \pair{\psi}{\UnitEff{\StSp}} \leq \pair{\psi}{f} \leq M\, \pair{\psi}{\UnitEff{\StSp}}$.
    If $t = \pair{\psi}{\UnitEff{\StSp}}$ were $0$, we would have $\pair{\psi}{f} = 0$ for
    all $f$ and hence $\psi = 0$. For $\psi \neq 0$ therefore $t > 0$, the functional
    $\psi / t$ is positive and normalised, and by Theorem~\ref{c02.thm.hoehe} it is a state
    $x$. If $\psi = t' x'$ is a second such representation, pairing with $\UnitEff{\StSp}$
    gives first $t' = t$ and then $x' = x$.
  \end{proof}

  So $\AffDualPos{\StSp}$ is the cone of the rays from the origin through $\StSp$: the set
  of unnormalised states. The second cone is described by Plávala as follows:

  \begin{lemma}[Positive functions and positive functionals]\label{c02.lem.dualkegel}\claim[proof-omitted]{c02.cl.dualkegel}{source-quoted}
    Let $f \in \Aff{\StSp}$. Then $f \in \AffPos{\StSp}$ if and only if
    $\pair{\psi}{f} \geq 0$ for all $\psi \in \AffDualPos{\StSp}$.
  \end{lemma}

  The cone of positive affine functions is thus the dual cone of the cone of unnormalised
  states, and the effects are the functions in it between $0$ and $\UnitEff{\StSp}$. At
  this place the framing of the course has the sentence
  % GAP-TJO: how TJO wants "elements of the dual space to states" (BRIEF.md:46) to be read
  % is not said in the brief (evidence/gaps.json, G-0010).  The sentence is therefore only
  % quoted; the mathematical statement after it is the chapter's own and not a reading of
  % his sentence.
  \begin{equation}\label{c02.eq.rahmen-dual}\claim{c02.cl.rahmen-dual}{author-stipulated}
    \begin{aligned}
      &\text{\tpcifenglish{Measurements/instruments/POVMs are then}{Messungen/Instrumente/POVMs sind dann}} \\
      &\text{\tpcifenglish{elements of the dual space to states.}{Elemente des Dualraums zu den Zuständen.}}
    \end{aligned}
  \end{equation}
  This sentence is quoted here and not interpreted. Separately from it, this block records
  as a mathematical statement: a single effect is an element of $\Aff{\StSp}$ with
  $0 \leq f \leq \UnitEff{\StSp}$, and $\Aff{\StSp}$ and the space $\AffDual{\StSp}$
  spanned by the states are dual to each other. A measurement is a list of such effects
  (\defref{D9}).

  \emph{In coordinates.} For $\StSp \subset \mathbb{R}^d$ as in
  Lemma~\ref{c02.lem.koordinaten} we write the state $\vb{r}$ as the vector
  $(1, \vb{r}) \in \mathbb{R}^{d+1}$ and the function $b + \vb{a}\cdot\vb{r}$ as the vector
  $(b, \vb{a})$. The pairing is then the ordinary scalar product of these two vectors,
  $b\cdot 1 + \vb{a}\cdot\vb{r}$. On the vectors $(1, \vb{r})$ effects are linear; on the
  points $\vb{r}$ themselves they are only affine. The two cones of the three theories are
  \begin{equation}\label{c02.eq.kegel-bit}\claim{c02.cl.kegel-bit}{derived}
    \begin{aligned}
      \AffDualPos{\Simplex{2}} &= \{ (t, R) : t \geq \abs{R} \}, \\
      \AffPos{\Simplex{2}} &= \{ (b, a) : b \geq \abs{a} \},
    \end{aligned}
  \end{equation}
  \begin{equation}\label{c02.eq.kegel-quadrat}\claim{c02.cl.kegel-quadrat}{derived}
    \begin{aligned}
      \AffDualPos{\StSp_{\mathrm{Q}}} &= \{ (t, R_x, R_y) : t \geq \max(\abs{R_x}, \abs{R_y}) \}, \\
      \AffPos{\StSp_{\mathrm{Q}}} &= \{ (b, a_x, a_y) : b \geq \abs{a_x} + \abs{a_y} \},
    \end{aligned}
  \end{equation}
  \begin{equation}\label{c02.eq.kegel-qubit}\claim{c02.cl.kegel-qubit}{derived}
    \begin{aligned}
      \AffDualPos{\StSp_{\mathrm{B}}} &= \{ (t, \vb{R}) : t \geq \norm{\vb{R}} \}, \\
      \AffPos{\StSp_{\mathrm{B}}} &= \{ (b, \vb{a}) : b \geq \norm{\vb{a}} \} ,
    \end{aligned}
  \end{equation}
  the bit here in the coordinate $r$ of equation~\eqref{c02.eq.bit-zentriert}.
  The first line is in each case the set of the vectors $t\,(1, \vb{r})$ with $t \geq 0$
  and $\vb{r} \in \StSp$, the second the set of the functions whose smallest value on
  $\StSp$ is non-negative; the smallest values were computed in Section~\ref{c02.s4}. For
  the bit and for the qubit the two sets are equal in these coordinates. For the square
  bit they differ: the vector $(1, 1, 1)$ lies in the first set and not in the second
  (Figure~\ref{c02.fig.kegel}). This is not a statement about other coordinates or other
  scalar products.

  \begin{tpcfigure}{c02.fig.kegel}
    \caption{The two cones of the square bit from \eqref{c02.eq.kegel-quadrat}. (a) The
      cone $\AffDualPos{\StSp}$ of unnormalised states. The state space $\StSp$ is its
      section with the plane $t = 1$; every ray from the origin through the cone meets
      $\StSp$ in exactly one state $x$. (b) The cone $\AffPos{\StSp}$ of positive affine
      functions; $b$ is plotted upwards. The effect set is the dashed double pyramid
      between $0$ and $\UnitEff{\StSp}$ of Figure~\ref{c02.fig.effektmengen}(b). The cone
      in (a) has a square cross-section with sides parallel to the axes, the one in (b) a
      cross-section turned by an eighth of a full turn and narrower: the two cones are not
      equal.}
  \end{tpcfigure}

  The unit effect plays a special role in $\Aff{\StSp}$: every affine function lies below
  a multiple of it, as used in the proof of Lemma~\ref{c02.lem.basis}. Such an element is
  called an order unit\zoomref{c02.f1}.
\end{vertiefung}

\section{Telling two preparations apart}\label{c02.s6}

Now to the task by which the three theories can be measured against each other. Given are
two states $x_1$ and $x_2$ of the same state space $\StSp$ and two numbers
$\lambda_1, \lambda_2 \geq 0$ with $\lambda_1 + \lambda_2 = 1$. Somebody prepares the state
$x_1$ with probability $\lambda_1$ and the state $x_2$ with probability $\lambda_2$ and
does not say which. We may ask one yes-no question $f$ and then have to guess. The
agreement is: on yes we guess $x_1$, on no we guess $x_2$.

We are right in two cases: $x_1$ was prepared and the answer is yes, or $x_2$ was prepared
and the answer is no. The success probability is therefore
\begin{equation}\label{c02.eq.erfolg}\claim{c02.cl.erfolg}{source-backed}
  P(f) = \lambda_1\, f(x_1) + \lambda_2\, \bigl(1 - f(x_2)\bigr) .
\end{equation}
We look for the largest value $P_{\mathrm{opt}}$ that $P(f)$ takes when $f$ runs through
all effects \srcref{plavala-2103.07469v2}{“Order unit norm, base norm and discrimination tasks”}.

\begin{tpcfigure}{c02.fig.aufgabe}
  \caption{The task. Left: $x_1$ or $x_2$ is prepared, then the question $f$ is asked and the guess follows the answer. Right: the probabilities of the four possible cases, in the upper row for the answer yes, in the lower row for no. In the two cases with the heavy frame the guess is right; their sum is $P(f)$ of equation~\eqref{c02.eq.erfolg}.}
\end{tpcfigure}

Two effects need no measurement. The zero function always answers no; we then always guess
$x_2$ and are right with probability $\lambda_2$. The unit effect always answers yes; we
always guess $x_1$ and are right with probability $\lambda_1$. Therefore in every theory
\begin{equation}\label{c02.eq.schranken}\claim{c02.cl.schranken}{derived}
  \max(\lambda_1, \lambda_2) \leq P_{\mathrm{opt}} \leq 1 .
\end{equation}

\begin{proposition}[When does the distinction succeed with certainty?]\label{c02.prop.perfekt}\claim[proof-included]{c02.cl.perfekt}{derived}
  Let $\lambda_1 > 0$ and $\lambda_2 > 0$. Then $P(f) = 1$ if and only if $f(x_1) = 1$ and
  $f(x_2) = 0$.
\end{proposition}
\begin{proof}
  From \eqref{c02.eq.erfolg} and $\lambda_1 + \lambda_2 = 1$ it follows that
  $1 - P(f) = \lambda_1\bigl(1 - f(x_1)\bigr) + \lambda_2\, f(x_2)$. Both terms are
  non-negative; the sum is $0$ only if both are $0$.
\end{proof}

In this chapter two states are called \emph{perfectly distinguishable} if there is an
effect with $f(x_1) = 1$ and $f(x_2) = 0$. The weights $\lambda_1$, $\lambda_2$ do not
matter for this as long as both are positive. If one of them is $0$, guessing alone
reaches the value $P = 1$, and the value says nothing about the two states.

On the set of effects the function $P$ is affine: in the coordinates $(b, \vb{a})$ of
Lemma~\ref{c02.lem.koordinaten},
\begin{equation}\label{c02.eq.erfolg-koordinaten}\claim{c02.cl.erfolg-koordinaten}{derived}
  P(f) = \lambda_2 + \delta\, b + \vb{v}\cdot\vb{a},
  \qquad \delta = \lambda_1 - \lambda_2, \quad \vb{v} = \lambda_1 \vb{r}_1 - \lambda_2 \vb{r}_2 ,
\end{equation}
where $\vb{r}_1$ and $\vb{r}_2$ are the coordinates of $x_1$ and $x_2$. If $\delta$
and $\vb{v}$ are not both $0$, then $P$ is not constant; the equation
$\delta\, b + \vb{v}\cdot\vb{a} = c$ then describes, for every number $c$, a hyperplane in
the coordinate space $\mathbb{R}^{d+1}$ of the effects, these hyperplanes are parallel, and
the level sets of $P$ are their intersections with the effect set. In the coordinates
$(b, a)$ of the bit they are lines, for the square bit planes. If $\delta = 0$ and
$\vb{v} = 0$, then $P$ is constant, and the whole effect set is a single level set. An
affine function takes its largest value
on the boundary, and in fact in a corner
(Figure~\ref{c02.fig.optimum})\zoomref{c02.v3}:

\begin{theorem}[An extreme effect attains the optimum]\label{c02.thm.extremal}\claim[proof-linked]{c02.cl.extremal}{derived}
  Let $\StSp$ be a state space, and let $x_1, x_2 \in \StSp$ and
  $\lambda_1, \lambda_2 \geq 0$ with $\lambda_1 + \lambda_2 = 1$ be given. Then there is an
  extreme effect $f$ with $P(f) = P_{\mathrm{opt}}$.
\end{theorem}

The theorem does not say that every optimal effect is extreme. In
Figure~\ref{c02.fig.optimum}(b) two extreme effects are optimal and with them all their
mixtures.

\begin{tpcfigure}{c02.fig.optimum}
  \caption{The success probability $P$ as a function on the effect set; the thin lines are
    level lines of $P$. (a) Classical bit with $x_1 = (\tfrac34, \tfrac14)$,
    $x_2 = (\tfrac14, \tfrac34)$ and $\lambda_1 = \lambda_2 = \tfrac12$: here
    $P = \tfrac12 + \tfrac14(q_0 - q_1)$, and the largest value $\tfrac34$ is taken only in
    the corner $(1,0)$. (b) Square bit with $x_1 = (1,1)$, $x_2 = (-1,-1)$ and
    $\lambda_1 = \lambda_2 = \tfrac12$: here $P = \tfrac12 + a_x + a_y$ independently of
    $b$; drawn is the section $b = \tfrac12$ through the double pyramid. The largest value
    $1$ is taken on the whole edge between $f_x^{+}$ and $f_y^{+}$. Optimal extreme effects
    are circled.}
\end{tpcfigure}

\begin{vertiefung}[id=c02.v3,
    titel={Why an extreme effect suffices},
    wozu={The core uses Theorem~\ref{c02.thm.extremal} without proof. Here it is proved
      for every state space.},
    voraussetzung={Blocks~\ref{c02.v1} and \ref{c02.v2}; continuous functions on closed,
      bounded sets. The core computes the optimum for the three theories directly in
      Proposition~\ref{c02.prop.optimum} and does not need this block.},
    zurueck={To tell two preparations apart as well as possible it suffices to go through
      the extreme effects.}]
  \begin{lemma}[The effect set is closed, bounded and convex]\label{c02.lem.kompakt}\claim[proof-included]{c02.cl.kompakt}{derived}
    Let $\StSp \subset \mathbb{R}^d$ be as in Lemma~\ref{c02.lem.koordinaten}. In the
    coordinates $(b, \vb{a}) \in \mathbb{R}^{d+1}$ the set $\EffAlg{\StSp}$ is closed,
    bounded and convex. It thus itself satisfies the conditions (S1) to (S5) of
    \defref{D3}.
  \end{lemma}
  \begin{proof}
    $\EffAlg{\StSp}$ is the set of all $(b, \vb{a})$ that satisfy the two inequalities
    $0 \leq b + \vb{a}\cdot\vb{r} \leq 1$ for every $\vb{r} \in \StSp$. Each single
    inequality describes a closed, convex set, and an intersection of such sets is again
    of this kind. $\EffAlg{\StSp}$ is bounded because, by the proof of
    Lemma~\ref{c02.lem.koordinaten}, $b$ and $\vb{a}$ are the solution of a system of
    linear equations whose right-hand sides $f(\vb{v}_i)$ lie in $[0,1]$.
  \end{proof}

  \begin{proof}[Proof of Theorem~\ref{c02.thm.extremal}]
    By equation~\eqref{c02.eq.erfolg-koordinaten} $P$ is a continuous function of the
    coordinates $(b, \vb{a})$. On the closed and bounded set $\EffAlg{\StSp}$ it takes its
    largest value in an effect $f^{*}$. By Lemma~\ref{c02.lem.kompakt},
    Theorem~\ref{c02.thm.zerlegung} also holds for $\EffAlg{\StSp}$: there are extreme
    effects $g_1, \dots, g_k$ and weights $\kappa_i \geq 0$ with sum $1$ and
    $f^{*} = \kappa_1 g_1 + \dots + \kappa_k g_k$. Since $P$ is affine, it follows as in
    equation~\eqref{c02.eq.mehrfach} that
    \begin{equation*}
      P_{\mathrm{opt}} = P(f^{*}) = \kappa_1 P(g_1) + \dots + \kappa_k P(g_k) .
    \end{equation*}
    No $P(g_i)$ is larger than $P_{\mathrm{opt}}$. A weighted mean of such numbers equals
    $P_{\mathrm{opt}}$ only if every $g_i$ with $\kappa_i > 0$ has the value
    $P(g_i) = P_{\mathrm{opt}}$. There is at least one such $g_i$.
  \end{proof}

  If $\StSp$ lies in an affine subspace of smaller dimension, the argument is applied in
  coordinates of that subspace.
\end{vertiefung}

For the three theories the optimum can be computed.

\begin{proposition}[The optimum in the three theories]\label{c02.prop.optimum}\claim[proof-included]{c02.cl.optimum}{derived}
  With $\delta = \lambda_1 - \lambda_2$ and $\vb{v} = \lambda_1 \vb{r}_1 - \lambda_2 \vb{r}_2$
  as in equation~\eqref{c02.eq.erfolg-koordinaten}:
  \begin{align}
    \text{\tpcifenglish{classical bit}{klassisches Bit}}\ (r = \pw_0 - \pw_1):\quad
      P_{\mathrm{opt}} &= \tfrac12\bigl(1 + \max(\abs{\delta},\, \abs{v})\bigr),
      \label{c02.eq.opt-bit} \\
    \text{\tpcifenglish{square bit}{Quadrat-Bit}}:\quad
      P_{\mathrm{opt}} &= \tfrac12\bigl(1 + \max(\abs{\delta},\, \abs{v_x},\, \abs{v_y})\bigr),
      \label{c02.eq.opt-quadrat} \\
    \text{Qubit}:\quad
      P_{\mathrm{opt}} &= \tfrac12\bigl(1 + \max(\abs{\delta},\, \norm{\vb{v}})\bigr).
      \label{c02.eq.opt-qubit}
  \end{align}
  For the classical bit, in the coordinates $x_1 = (\pw_0, \pw_1)$,
  $x_2 = (\pw'_0, \pw'_1)$, this is the same as
  \begin{equation}\label{c02.eq.opt-bit-p}
    P_{\mathrm{opt}} = \max(\lambda_1 \pw_0,\, \lambda_2 \pw'_0) + \max(\lambda_1 \pw_1,\, \lambda_2 \pw'_1) .
  \end{equation}
\end{proposition}
\begin{proof}
  In equation~\eqref{c02.eq.erfolg-koordinaten} we first hold $b$ fixed and choose $\vb{a}$
  as well as possible. Allowed are all $\vb{a}$ whose length is at most
  $m = \min(b, 1-b)$, where the length means $\abs{a}$, $\abs{a_x} + \abs{a_y}$ or
  $\norm{\vb{a}}$ according to the theory. The largest value of
  $\vb{v}\cdot\vb{a}$ is then $m\, L$, with $L = \abs{v}$ for the bit,
  $L = \max(\abs{v_x}, \abs{v_y})$ for the square bit and $L = \norm{\vb{v}}$ for the qubit.
  For the bit one chooses $a = \pm m$ with the sign of $v$. For the square bit one puts the
  whole of $m$ on the coordinate in which $\vb{v}$ has the larger modulus. For the qubit
  one chooses $\vb{a}$ in the direction of $\vb{v}$.

  It remains to maximise $\lambda_2 + \delta\, b + L \min(b, 1-b)$ over $b \in [0,1]$. On
  $[0, \tfrac12]$ and on $[\tfrac12, 1]$ this expression is a straight line in $b$; its
  largest value is therefore at $b = 0$, $b = \tfrac12$ or $b = 1$. The three values are
  $\lambda_2$, $\lambda_2 + \tfrac12\delta + \tfrac12 L = \tfrac12(1 + L)$ and
  $\lambda_2 + \delta = \lambda_1$. Because
  $\max(\lambda_1, \lambda_2) = \tfrac12(1 + \abs{\delta})$, the largest of the three is
  $\tfrac12(1 + \max(\abs{\delta}, L))$.

  On \eqref{c02.eq.opt-bit-p}: inserting \eqref{c02.eq.bit-effekt} into
  \eqref{c02.eq.erfolg} gives $P = \lambda_2 + d_0 q_0 + d_1 q_1$ with
  $d_0 = \lambda_1 \pw_0 - \lambda_2 \pw'_0$ and $d_1 = \lambda_1 \pw_1 - \lambda_2 \pw'_1$.
  This is largest if one chooses $q_i = 1$ where $d_i > 0$, and $q_i = 0$ otherwise.
  That gives $\lambda_2 \pw'_0 + \max(d_0, 0) + \lambda_2 \pw'_1 + \max(d_1, 0)$, and this is
  the right-hand side of \eqref{c02.eq.opt-bit-p}.
\end{proof}

The proof supplies the optimal effects as well. If $\abs{\delta}$ is the largest of the
numbers in the bracket, guessing is optimal: the unit effect if $\lambda_1$ is the larger
weight, otherwise the zero function. Otherwise an extreme effect with $b = \tfrac12$ is
optimal: for the square bit one of the four effects $f_x^{\pm}$, $f_y^{\pm}$, for the qubit
the effect $f_{\vb{n}}$ with $\vb{n} = \vb{v}/\norm{\vb{v}}$. The three formulas differ
only in how the length of $\vb{v}$ is measured (Figure~\ref{c02.fig.laengen})\zoomref{c02.v4}.

\begin{tpcfigure}{c02.fig.laengen}
  \caption{How well can a state $x_1$ be told apart from the centre $x_2$ of the state
    space if both are prepared equally often? By Proposition~\ref{c02.prop.optimum} we
    then have $P_{\mathrm{opt}} = \tfrac12 + \tfrac14 L$, where $L$ is the length of the
    position vector of $x_1$ in the respective theory. Marked are the positions of $x_1$
    with $P_{\mathrm{opt}} = \tfrac58$ and $P_{\mathrm{opt}} = \tfrac34$: (a) points,
    (b) squares, (c) circles. More than $\tfrac34$ cannot be reached against the centre in
    any of the three theories.}
\end{tpcfigure}

\begin{vertiefung}[id=c02.v4,
    titel={The qubit optimum as a trace norm},
    wozu={Equation~\eqref{c02.eq.opt-qubit} can be written without coordinates, directly
      with the density matrices of the two states.},
    voraussetzung={Section~\ref{c02.s4.3} and Proposition~\ref{c02.prop.optimum}.
      The core does not need this block.},
    zurueck={The length $\max(\abs{\delta}, \norm{\vb{v}})$ in
      equation~\eqref{c02.eq.opt-qubit} is a property of the matrix
      $\lambda_1\rho_1 - \lambda_2\rho_2$ alone.}]
  For a Hermitian $2\times 2$ matrix $\op{D}$ with the eigenvalues $\kappa_1$ and $\kappa_2$ we
  call $\norm{\op{D}}_1 = \abs{\kappa_1} + \abs{\kappa_2}$ its \emph{trace norm}.

  \begin{proposition}[Trace-norm formula]\label{c02.prop.spurnorm}\claim[proof-included]{c02.cl.spurnorm}{derived}
    Let $\rho_1 = \rho(\vb{r}_1)$ and $\rho_2 = \rho(\vb{r}_2)$ be two states of the qubit
    and $\lambda_1, \lambda_2 \geq 0$ with $\lambda_1 + \lambda_2 = 1$. Then
    \begin{equation}\label{c02.eq.spurnorm}
      P_{\mathrm{opt}} = \tfrac12\bigl(1 + \norm{\lambda_1 \rho_1 - \lambda_2 \rho_2}_1\bigr) .
    \end{equation}
  \end{proposition}
  \begin{proof}
    By equation~\eqref{c02.eq.bloch},
    $\op{D} = \lambda_1\rho_1 - \lambda_2\rho_2
       = \tfrac12\bigl(\delta\, \Id + v_x \pauli{x} + v_y \pauli{y} + v_z \pauli{z}\bigr)$.
    The calculation of Proposition~\ref{c02.prop.bloch}(2) gives for $\op{D}$ the eigenvalues
    $\tfrac12(\delta + \norm{\vb{v}})$ and $\tfrac12(\delta - \norm{\vb{v}})$. With
    $L = \norm{\vb{v}}$ we therefore have
    $\norm{\op{D}}_1 = \tfrac12\bigl(\abs{\delta + L} + \abs{\delta - L}\bigr)$. For
    $\delta \geq L$ this is $\delta$, for $\delta \leq -L$ it is $-\delta$, and for
    $-L \leq \delta \leq L$ it is $L$; in every case $\max(\abs{\delta}, L)$.
    Inserting into \eqref{c02.eq.opt-qubit} gives the claim.
  \end{proof}

  Plávala proves a formula of the same shape for every state space. In place of the trace
  norm it has a norm on $\AffDual{\StSp}$, which is not introduced here
  \srcref{plavala-2103.07469v2}{“Order unit norm, base norm and discrimination tasks”}.
\end{vertiefung}

What follows for pure states?

\begin{corollary}[Pure states are not always perfectly distinguishable]\label{c02.cor.rein}\claim[proof-included]{c02.cl.rein}{derived}
  \begin{enumerate}
    \item In the classical bit and in the square bit any two different pure states are
      perfectly distinguishable.
    \item Two states of the qubit are perfectly distinguishable if and only if both are
      pure and opposite to each other: $\vb{r}_2 = -\vb{r}_1$ with $\norm{\vb{r}_1} = 1$.
    \item The pure states $x_1 = (1,0,0)$ and $x_2 = (0,1,0)$ of the qubit are different
      and not perfectly distinguishable. For $\lambda_1 = \lambda_2 = \tfrac12$,
      \begin{equation}\label{c02.eq.q1}
        P_{\mathrm{opt}} = \tfrac12\Bigl(1 + \tfrac12\sqrt{2}\Bigr) = \tfrac14\bigl(2 + \sqrt{2}\bigr) < 1 .
      \end{equation}
  \end{enumerate}
\end{corollary}
\begin{proof}
  (1) For the bit the effect $(1,0)$ does what is required for $x_1 = s_0$, $x_2 = s_1$.
  Two different corners of the square differ in at least one coordinate; it is $1$ in one
  corner and $-1$ in the other. The corresponding effect $f_x^{\pm}$ or $f_y^{\pm}$ then
  has the values $1$ and $0$.

  (2) By Proposition~\ref{c02.prop.perfekt} we may assume
  $\lambda_1 = \lambda_2 = \tfrac12$. Then $\delta = 0$ and
  $\vb{v} = \tfrac12(\vb{r}_1 - \vb{r}_2)$, and equation~\eqref{c02.eq.opt-qubit} gives
  $P_{\mathrm{opt}} = 1$ if and only if $\norm{\vb{r}_1 - \vb{r}_2} = 2$. Because
  $\norm{\vb{r}_1 - \vb{r}_2}^2 = \norm{\vb{r}_1}^2 + \norm{\vb{r}_2}^2 - 2\,\vb{r}_1\cdot\vb{r}_2 \leq 4$
  this is possible only if both lengths are $1$ and $\vb{r}_1\cdot\vb{r}_2 = -1$, that is,
  $\vb{r}_2 = -\vb{r}_1$.

  (3) Here $\vb{v} = \tfrac12(1, -1, 0)$ and $\norm{\vb{v}} = \tfrac12\sqrt{2}$.
\end{proof}

The numerical value in \eqref{c02.eq.q1} is approximately \num[c02.cl.level-q1]{0.8536}{}.
Part (2) holds for the ball. It cannot be carried over to other state spaces; in the square
even neighbouring corners are perfectly distinguishable\zoomref{c02.v5}.

The opposite case occurs too: the states are different, and still guessing is best. For
the qubit states $x_1 = (0, 0, \tfrac14)$ and $x_2 = (0, 0, -\tfrac14)$ with the weights
$\lambda_1 = \tfrac{9}{10}$, $\lambda_2 = \tfrac{1}{10}$ we have $\delta = \tfrac45$ and
$\norm{\vb{v}} = \tfrac14$. Equation~\eqref{c02.eq.opt-qubit} gives
$P_{\mathrm{opt}} = \num[c02.cl.level-q2]{\tfrac{9}{10}}{}$, and the unit effect already
reaches this: always guess $x_1$.

\begin{vertiefung}[id=c02.v5,
    titel={Polygons and the disc},
    wozu={Between the square and the ball lies a whole family of state spaces in the
      plane. It shows how the distinguishability of pure states depends on the shape of
      the state space, and what can be said about the relation of the two cones.},
    voraussetzung={Core up to Corollary~\ref{c02.cor.rein}; for the last part
      block~\ref{c02.v2}. The core does not need this block.},
    zurueck={Whether two pure states can be told apart with certainty depends on the shape
      of the state space and not only on the two points.}]
  For $n \geq 3$ let $\StSp_n$ be the regular $n$-gon with circumradius $1$ and the corners
  \begin{equation}\label{c02.eq.vieleck}\claim{c02.cl.vieleck}{derived}
    \vb{v}_k = \bigl(\cos\tfrac{2\pi k}{n},\ \sin\tfrac{2\pi k}{n}\bigr), \quad k = 0, \dots, n-1,
    \qquad c_n = \cos\tfrac{\pi}{n} ,
  \end{equation}
  and $\StSp_{\mathrm{D}} = \{ \vb{r} \in \mathbb{R}^2 : \norm{\vb{r}} \leq 1 \}$ the unit
  disc. The pure states of $\StSp_n$ are the corners $\vb{v}_k$, those of
  $\StSp_{\mathrm{D}}$ the points of the boundary circle. $\StSp_3$ is a triangle and hence
  a simplex, the state space of a classical theory with three pure states. $\StSp_4$ is a
  square whose corners lie on the axes; it arises from $\StSp_{\mathrm{Q}}$ by a rotation
  and a reduction in size.

  \begin{proposition}[Extreme effects of the polygons and of the disc]\label{c02.prop.vieleck}\claim[proof-omitted]{c02.cl.vieleck-effekte}{derived}
    Let $\vb{m}_j = \bigl(\cos\tfrac{(2j+1)\pi}{n},\ \sin\tfrac{(2j+1)\pi}{n}\bigr)$ for
    $j = 0, \dots, n-1$, and let
    \begin{equation}\label{c02.eq.vieleck-effekte}
      f_k(\vb{r}) = \frac{c_n + \vb{v}_k\cdot\vb{r}}{1 + c_n},
      \qquad
      g_j(\vb{r}) = \frac12 + \frac{\vb{m}_j\cdot\vb{r}}{2 c_n} .
    \end{equation}
    The extreme effects of $\StSp_n$ are $0$ and $\UnitEff{\StSp}$ and in addition: for
    odd $n$ the effects $f_k$ and $\UnitEff{\StSp} - f_k$ with $k = 0, \dots, n-1$; for
    even $n$ the effects $g_j$ with $j = 0, \dots, n-1$.
    For odd $n$ these are $2n + 2$ extreme effects, for even $n$ there are $n + 2$.
    The extreme effects of $\StSp_{\mathrm{D}}$ are $0$, $\UnitEff{\StSp}$ and
    $\tfrac12(1 + \vb{n}\cdot\vb{r})$ for every unit vector $\vb{n}$ of the plane.
  \end{proposition}

  For odd $n$ the effect $f_k$ has the value $1$ in the corner $\vb{v}_k$ and the value $0$
  on the opposite edge. For even $n$ the effect $g_j$ has the value $1$ on one edge and the
  value $0$ on the opposite one. For $n = 4$ the $g_j$ are the four effects $f_x^{\pm}$,
  $f_y^{\pm}$ of the square in rotated coordinates. The derivation is kept with the
  evidence of the chapter; for $n = 3, \dots, 8$ the lists were in addition checked by
  enumerating all corners of the effect set.

  \emph{Distinguishing.} The formula for the disc is the same as for the ball,
  equation~\eqref{c02.eq.opt-qubit} with vectors of the plane. For the polygons one goes
  through the finitely many extreme effects (Theorem~\ref{c02.thm.extremal}). Three
  results, all for $\lambda_1 = \lambda_2 = \tfrac12$:
  \begin{enumerate}
    \item In the pentagon the neighbouring pure states $\vb{v}_0$ and $\vb{v}_1$ are not
      perfectly distinguishable: the optimum is
      $P_{\mathrm{opt}} = c_5 = \tfrac14(1 + \sqrt{5})$, approximately
      \num[c02.cl.level-p1]{0.8090}{}.
    \item In the hexagon $\vb{v}_0 = (1,0)$ and
      $\vb{v}_2 = (-\tfrac12, \tfrac12\sqrt{3})$ are pure states and perfectly
      distinguishable, $P_{\mathrm{opt}} = \num[c02.cl.level-h1]{1}{}$, although they are
      not opposite to each other. The effect $g_5$ has the value $1$ in $\vb{v}_0$ and the
      value $0$ in $\vb{v}_2$.
    \item The same pair of points $x_1 = (1,0)$, $x_2 = (-\tfrac34, \tfrac14)$ lies in the
      square $\StSp_4$, in the pentagon $\StSp_5$ and in the disc $\StSp_{\mathrm{D}}$
      (Figure~\ref{c02.fig.vielecke}). The optima are
      \begin{equation}\label{c02.eq.gleiches-paar}\claim{c02.cl.gleiches-paar}{derived}
        P_{\mathrm{opt}}(\StSp_4) = 1, \qquad
        P_{\mathrm{opt}}(\StSp_5) = \frac{55 - 7\sqrt{5}}{40}, \qquad
        P_{\mathrm{opt}}(\StSp_{\mathrm{D}}) = \frac{8 + 5\sqrt{2}}{16} ,
      \end{equation}
      approximately \num[c02.cl.level-c1]{1}{}, \num[c02.cl.level-c2]{0.9837}{} and
      \num[c02.cl.level-c3]{0.9419}{}.
  \end{enumerate}
  The third statement depends on how the three state spaces lie relative to each other.
  With the axis-parallel square $\StSp_{\mathrm{Q}} = [-1,1]^2$ and the unit disc in the
  same coordinates there is no such pair: every pair of points of the disc that is
  perfectly distinguishable in $\StSp_{\mathrm{Q}}$ consists of two opposite boundary
  points of the disc and is perfectly distinguishable there as well.

  \begin{tpcfigure}{c02.fig.vielecke}
    \caption{The same pair of points $x_1 = (1,0)$, $x_2 = (-\tfrac34, \tfrac14)$ in three
      state spaces, each with the level lines of an optimal extreme effect. (a) In the
      square $\StSp_4$ the point $x_2$ lies on the edge on which the effect has the value
      $0$; the distinction succeeds with certainty. (b) In the pentagon $x_2$ lies in the
      interior, just beside the edge with $f = 0$. (c) In the disc the line $f = 0$
      touches the boundary in a single point, and $x_2$ lies clearly beside it. The optima
      are in \eqref{c02.eq.gleiches-paar}.}
  \end{tpcfigure}

  \emph{The two cones.} In the coordinates of block~\ref{c02.v2} the $n$-gon with
  circumradius $\gamma > 0$ has the state cone
  $C_{n,\gamma} = \{ t\,(1, \gamma\,\vb{r}) : t \geq 0,\ \vb{r} \in \StSp_n \}$.

  \begin{proposition}[When the cones of the polygons are equal]\label{c02.prop.selbstdual}\claim[proof-omitted]{c02.cl.selbstdual}{derived}
    Let $C_{n,\gamma}^{*} = \{ (s, \vb{a}) : s\,t + \vb{a}\cdot\vb{R} \geq 0
    \ \text{\tpcifenglish{for all}{für alle}}\ (t, \vb{R}) \in C_{n,\gamma} \}$ be the dual
    cone with respect to the ordinary scalar product of $\mathbb{R}^3$.
    \begin{enumerate}
      \item For $\gamma = 1$ we have $C_{n,1}^{*} \neq C_{n,1}$ for every $n \geq 3$.
      \item For odd $n$ we have $C_{n,\gamma}^{*} = C_{n,\gamma}$ if and only if
        $\gamma^2 = 1/c_n$.
      \item For even $n$ we have $C_{n,\gamma}^{*} \neq C_{n,\gamma}$ for every
        $\gamma > 0$.
      \item For the disc the cone $\{ (t, \vb{R}) : t \geq \norm{\vb{R}} \}$ is equal to
        its dual cone.
    \end{enumerate}
  \end{proposition}

  Statements (1) to (3) were in addition checked numerically for $n = 3, \dots, 8$. They
  concern the equality of two sets in fixed coordinates with the ordinary scalar product.
  For every $n$ and every $\gamma$ there is an invertible linear map that maps
  $C_{n,\gamma}^{*}$ onto $C_{n,\gamma}$, a rotation combined with a scaling; that is a
  weaker statement than equality. Whether for even $n$ another scalar product can be found
  for which the cones become equal is not examined here. The triangle too falls under (1):
  in the probability coordinates $(\pw_1, \pw_2, \pw_3)$ the state cone of the simplex is
  equal to its dual cone, in the coordinates of the regular triangle with circumradius $1$
  it is not. The two coordinate systems are not related by a rotation.
\end{vertiefung}

\section{The game “Distinguish”}\label{c02.s7}

A game belongs to this chapter. This page is its manual.

\emph{Rules.} A level specifies a theory, two states $x_1$, $x_2$ and two weights
$\lambda_1$, $\lambda_2$. One plays the task of Section~\ref{c02.s6}: an effect $f$ is
dragged to a place in the effect set. The measurement is the pair $f$,
$\UnitEff{\StSp} - f$; on yes the guess is $x_1$, on no it is $x_2$. There is no chance and
no time limit.

\emph{The success quantity} is $P(f)$ of equation~\eqref{c02.eq.erfolg}, the probability
of guessing right. It is the only physical quantity of the game. The game displays it for
the chosen effect.

\emph{The reward} is computed from $P(f)$ and $P_{\mathrm{opt}}$ and is a decision of game
design, not physics. Three stars are given only for $P(f) = P_{\mathrm{opt}}$; the
comparison is exact, without tolerance. Below that there are two tiers with fixed
thresholds $0 < \varepsilon_2 < \varepsilon_1$: two stars for
$P_{\mathrm{opt}} - P(f) \leq \varepsilon_2$, otherwise one star for
$P_{\mathrm{opt}} - P(f) \leq \varepsilon_1$, otherwise none. The two thresholds are
chosen design values, not derived and not validated; their values are in the rules panel
of the game and not in this book.

\emph{The proven optimum} $P_{\mathrm{opt}}$ of every level is in
Table~\ref{c02.tab.level}. It follows from Proposition~\ref{c02.prop.optimum} and, by
Theorem~\ref{c02.thm.extremal}, is attained by an extreme effect. More is not possible in
any level.

\begin{tpctable}{c02.tab.level}
  \caption{The levels of the first release. Bit: states as $(\pw_0, \pw_1)$, effects as
    $(q_0, q_1)$. Square bit and qubit: states as $\vb{r}$. Last column: an optimal
    extreme effect; in Q1, $\vb{n} = (1, -1, 0)/\sqrt{2}$.}
  \begin{tabular}{llllll}
    \toprule
    Level & $x_1$ & $x_2$ & $(\lambda_1, \lambda_2)$ & $P_{\mathrm{opt}}$ & optimal \\
    \midrule
    B1 & $(1, 0)$ & $(0, 1)$ & $(\tfrac12, \tfrac12)$ & \num[c02.cl.level-b1]{1}{} & $(1, 0)$ \\
    B2 & $(\tfrac34, \tfrac14)$ & $(\tfrac14, \tfrac34)$ & $(\tfrac12, \tfrac12)$ & \num[c02.cl.level-b2]{\tfrac34}{} & $(1, 0)$ \\
    S1 & $(1, 0)$ & $(-1, 0)$ & $(\tfrac12, \tfrac12)$ & \num[c02.cl.level-s1]{1}{} & $f_x^{+}$ \\
    S2 & $(1, 1)$ & $(-1, -1)$ & $(\tfrac12, \tfrac12)$ & \num[c02.cl.level-s2]{1}{} & $f_x^{+}$, $f_y^{+}$ \\
    Q1 & $(1, 0, 0)$ & $(0, 1, 0)$ & $(\tfrac12, \tfrac12)$ & \num[c02.cl.level-q1]{\tfrac14(2 + \sqrt{2})}{} & $f_{\vb{n}}$ \\
    Q2 & $(0, 0, \tfrac14)$ & $(0, 0, -\tfrac14)$ & $(\tfrac{9}{10}, \tfrac{1}{10})$ & \num[c02.cl.level-q2]{\tfrac{9}{10}}{} & $\UnitEff{\StSp}$ \\
    \bottomrule
  \end{tabular}
\end{tpctable}

\begin{tpcfigure}{c02.fig.unterscheiden}
  \demo{c02.demo.distinguish}
  \caption{One level in each of the three theories with the level lines of an optimal
    extreme effect (dashed). (a) B2: the effect $(1,0)$ has the value $\tfrac34$ in $x_1$
    and the value $\tfrac14$ in $x_2$. (b) S2: the effect $f_x^{+}$ separates the two
    corners with certainty; so does $f_y^{+}$. (c) Q1, section $r_z = 0$ through the qubit
    state ball: no family of parallel lines has $x_1$ on $f = 1$ and at the same time $x_2$ on
    $f = 0$.}
\end{tpcfigure}

In S1 both states are mixed and nevertheless perfectly distinguishable. In S2 two extreme
effects are optimal and with them the whole edge between them
(Figure~\ref{c02.fig.optimum}(b)). In Q2 guessing is optimal: the best effect is the unit
effect, always guess $x_1$.

\section{From here to quantum theory}\label{c02.s8}

The framework of this chapter singles out none of the three state spaces: segment, square
and ball all satisfy the conditions of \defref{D3}. Which state space belongs to a given
system the framework does not say. In quantum theory the state spaces of a particular
family occur, to which the ball of Section~\ref{c02.s4.3} belongs as the smallest
nontrivial example. For the step to it the framing of the course has the sentence:
% GAP-TJO: how TJO wants the step from GPTs to quantum theory named is open
% (evidence/gaps.json, G-0012).  The label "choice of model" of docs/plan/PLAN.md,
% section 13.1, A11 is not his wording and therefore does not stand in the text; his
% sentence is the following one (BRIEF.md:47).
\begin{equation}\label{c02.eq.rahmen-spezialisierung}\claim{c02.cl.rahmen-spezialisierung}{author-stipulated}
  \begin{aligned}
    &\text{\tpcifenglish{The GPT presentation allows us to present}{Die GPT-Darstellung erlaubt uns, die Postulate}} \\
    &\text{\tpcifenglish{the postulates of a physical theory which we}{einer physikalischen Theorie anzugeben, die wir}} \\
    &\text{\tpcifenglish{subsequently specialise to QM next.}{anschließend zur Quantenmechanik spezialisieren.}}
  \end{aligned}
\end{equation}

\begin{fundament}[id=c02.f1,
    titel={Order unit spaces},
    wozu={The effects of a state space form an interval in an ordered vector space with
      an order unit. Here this structure is stated as Plávala gives it for finite
      dimension, and it is said what the course still owes beyond that.},
    voraussetzung={Block~\ref{c02.v2}; partial orders. Neither the core nor the other
      blocks need this block.},
    zurueck={A theory can be started from its states or from its effects; in finite
      dimension both lead to the same framework.}]
  On $\Aff{\StSp}$ there is a natural order: $f \leq g$ means $f(x) \leq g(x)$ for all
  $x \in \StSp$, that is, $g - f \in \AffPos{\StSp}$. For Plávala an \emph{ordered vector
  space} is a vector space with a relation $\leq$ that is reflexive, antisymmetric and
  transitive and compatible with addition and with multiplication by non-negative numbers
  \srcref{plavala-2103.07469v2}{appendix “Convex cones and ordered vector spaces”}.

  \tpcdef{D6}

  \begin{proposition}[Affine functions as an order unit space]\label{c02.prop.oeraum}\claim[proof-omitted]{c02.cl.oeraum}{source-backed}
    For every state space $\StSp$, $(\Aff{\StSp}, \leq, \UnitEff{\StSp})$ is an order unit
    space, and the effects are in it the order interval between $0$ and the order unit:
    \begin{equation}\label{c02.eq.intervall}
      \EffAlg{\StSp} = \{ f \in \Aff{\StSp} : 0 \leq f \leq \UnitEff{\StSp} \} .
    \end{equation}
  \end{proposition}

  Plávala proves the three ingredients one by one: that $\Aff{\StSp}$ with this order is an
  ordered vector space, that $\UnitEff{\StSp}$ is an order unit, and
  equation~\eqref{c02.eq.intervall}. Source:
  \srcref{plavala-2103.07469v2}{“Effect algebra”}.
  Figure~\ref{c02.fig.ordnungsintervall} shows the interval for the classical bit.

  \begin{tpcfigure}{c02.fig.ordnungsintervall}
    \caption{The effect set of the classical bit as an order interval, in the coordinates
      $(q_0, q_1)$ of Proposition~\ref{c02.prop.bit}. The positive cone $\AffPos{\StSp}$ is
      the quadrant to the upper right of the origin. The condition
      $f \leq \UnitEff{\StSp}$ says that $\UnitEff{\StSp} - f$ lies in this cone; that is
      the quadrant to the lower left of $\UnitEff{\StSp}$. The intersection of the two
      quadrants is the unit square $\EffAlg{\StSp}$.}
  \end{tpcfigure}

  The way also leads in the other direction, from the effects to the states. A cone $C$ in
  a vector space $V$ is called \emph{pointed} if no vector other than $0$ lies in $C$
  together with its negative, and \emph{generating} if $C$ spans the space $V$.

  \begin{proposition}[From the effects back to the states]\label{c02.prop.rueckweg}\claim[proof-omitted]{c02.cl.rueckweg}{source-backed}
    Let $V$ be a real, finite-dimensional vector space, $C \subset V$ a convex, closed,
    pointed and generating cone and $\OrdUnit \in C$ an order unit. Let
    $E = \{ v \in V : 0 \leq v \leq \OrdUnit \}$ and
    \begin{equation}\label{c02.eq.rueckweg}
      S(E) = \{ \psi \in C^{*} : \pair{\psi}{\OrdUnit} = 1 \} ,
    \end{equation}
    where $C^{*}$ is the dual cone of $C$. Then $S(E)$ is a state space, and its effect set
    is again $E$: we have $\EffAlg{S(E)} = E$.
  \end{proposition}

  Plávala indicates the proof and concludes that the framework that starts from such
  intervals is the same as the one that starts from state spaces
  \srcref{plavala-2103.07469v2}{“Connection to abstract convex effect algebras and order unit spaces”}.
  The abstract notion of an effect algebra, which this chapter does not need, is found
  there as well.

  \emph{What remains open.} At this place the framing of the course has the sentence
  % GAP-TJO: BRIEF.md:47 requires "Archimedean order unit spaces".  The word "Archimedean"
  % does not occur in the source (Plávala v2), and no source on disk defines it
  % (evidence/gaps.json, G-0002).  What it means for the course, how it relates to
  % definition D6 and how the sentence is to be read in infinite dimension is for TJO to
  % say.
  \begin{equation}\label{c02.eq.rahmen-archimedisch}\claim{c02.cl.rahmen-archimedisch}{author-stipulated}
    \begin{aligned}
      &\text{\tpcifenglish{All physical theories, for me,}{Alle physikalischen Theorien sind für mich}} \\
      &\text{\tpcifenglish{are Archimedean order unit spaces.}{archimedische Ordnungseinheitsräume.}}
    \end{aligned}
  \end{equation}
  This chapter covers only part of it. First, \defref{D6} is the definition of the source,
  and the property “Archimedean” does not occur in it; what it means and why the course
  requires it is not said here. Second, all statements of this chapter are proved or
  quoted under condition (S5), that is, for finite-dimensional state spaces. About
  infinite-dimensional theories it says nothing.
\end{fundament}
